\section{模型概率、置信度与校准}

课程问答指出，输出 token 的平均 log probability 与正确性有相关性，但通常校准不足：错误回答也可能得到很高内部概率，且模型在被纠正时仍坚持原答案。

可用可靠性图或 Expected Calibration Error 衡量：

$$
\operatorname{ECE}=\sum_{b=1}^{B}\frac{|S_b|}{N}
\left|\operatorname{acc}(S_b)-\operatorname{conf}(S_b)\right|.
$$

\begin{itemize}
    \item $S_b$：预测置信度落入第 $b$ 个区间的样本；
    \item $\operatorname{acc}(S_b)$：该区间真实正确率；
    \item $\operatorname{conf}(S_b)$：该区间平均置信度；
    \item ECE 越小，置信度越接近经验正确率。
\end{itemize}

\begin{knowledgebox}{选择器最好学习可验证信号，而非只读模型自信}
代码任务可加入编译、样例、变异测试、复杂度分析和行为聚类；研究 agent 可加入来源质量、引用一致性和交叉证据。外部信号通常比生成模型的自报置信度更可靠。
\end{knowledgebox}

\subsection{本章小结}

语言模型概率不是经过校准的“答案正确率”。候选排序、停止条件和是否触发搜索都应结合外部验证与专门校准，而不能只看原始 token probability。

